\chapter{Introduction}
\label{ch:introduction}

\section{Purpose and Scope}
\label{sec:purpose_scope}
This report presents the SRR-level architecture baseline of the LINKU mission and defines the corresponding mission and system requirements. It addresses the design of the LINKU multi-satellite system, including the carrier--payload configuration, the SAT-BC release from SAT-A, the SAT-B/SAT-C separation and tether deployment sequence, the communications architecture, the attitude and stabilization functions, and the subsystem allocation required to support the tether experiment.

The document reviews the LINKU Concept of Operations (CONOPS), with emphasis on the transition from the integrated SAT-ABC configuration to the coupled SAT-BC configuration and finally to the deployed SAT-B/SAT-C tethered configuration. This operational sequence provides the basis for the payload design, subsystem analyses, and derived requirements presented in the following chapters.

The report consolidates results obtained in previous LINKU technical studies, including tether modelling, tether-material and environmental considerations, and tethered-system dynamics. These results are used to define a consistent SRR-level requirements baseline, while detailed theoretical developments and simulations are referenced where relevant but are not reproduced in full.

This document supports the Stage-1 development of the LINKU Engineering Model. It does not constitute a final flight-model design or a flight qualification document.

\section{Document Structure and Reference Documents}
\label{sec:document_structure_references}
The document is organized to progress from mission definition to system requirements and verification. Chapter~\ref{ch:mission_overview_conops} presents the mission background, objectives, scenario, and Concept of Operations. Chapter~\ref{ch:system_configuration} describes the physical configuration of SAT-A, SAT-B, SAT-C, and the coupled SAT-BC assembly. Chapter~\ref{ch:payload_tether_experiment} describes the tether experiment, including the SAT-B/SAT-C separation and tether deployment mechanism, tether storage and braking concept, and measurement instrumentation. Chapter~\ref{ch:satellite_subsystems} presents the satellite subsystem analyses, including the electrical power system, communications system, attitude determination and control system, and structural configuration. Chapters~\ref{ch:mission_requirements} and~\ref{ch:system_requirements} define the mission-level and system-level requirements, and Chapter~\ref{ch:verification_approach} introduces the verification approach for the SRR-level baseline.

The main reference documents used for this report are the previous LINKU technical deliverables: D2, focused on tether modelling; D3, focused on tether materials and environmental considerations; and D4, focused on the dynamic analysis of the tethered satellite system. These internal references are complemented by the LINKU feasibility and project-management documents, applicable system-engineering guidance, CubeSat design considerations, and subsystem-level references used in the corresponding technical sections.

